\chapter{Machine learning for clinical prediction}
\label{ch:ml-clinical}

\begin{quote}
\emph{``A model that works perfectly on your training data and fails on the next patient is not a model --- it is a memory.''}
\end{quote}

% ============================================================
\section{When to use ML vs traditional statistics}

Machine learning and traditional statistics answer different questions:

\begin{itemize}
  \item \textbf{Traditional statistics} (regression, t-tests): ``Which risk factors are associated with heart disease, and by how much?'' The goal is \emph{inference} --- understanding relationships.
  \item \textbf{Machine learning}: ``Given this patient's data, what is their probability of heart disease?'' The goal is \emph{prediction} --- accuracy on new patients.
\end{itemize}

\begin{clinicalnote}
Use logistic regression when you need to explain \emph{why} (odds ratios, confidence intervals, p-values for each risk factor). Use ML when you need the best possible \emph{prediction} and interpretability is secondary. In practice, many clinical studies use both: regression to understand the biology, ML to build the screening tool.
\end{clinicalnote}

\subsection{When NOT to use ML}

\begin{itemize}
  \item Small datasets ($< 200$ patients): ML overfits easily. Stick with logistic regression.
  \item When you need regulatory approval: the FDA requires interpretable models for most clinical decision support tools.
  \item When a simple rule works: if ``glucose $> 126$'' classifies 90\% of diabetics correctly, you do not need a random forest.
\end{itemize}

\begin{minted}{python}
import numpy as np
import pandas as pd
import matplotlib.pyplot as plt
import seaborn as sns
from sklearn.model_selection import train_test_split, cross_val_score
from sklearn.preprocessing import StandardScaler
from sklearn.linear_model import LogisticRegression
from sklearn.tree import DecisionTreeClassifier, plot_tree
from sklearn.ensemble import RandomForestClassifier
from sklearn.metrics import (classification_report, confusion_matrix,
                             roc_curve, roc_auc_score, RocCurveDisplay)
\end{minted}

% ============================================================
\section{Loading clinical data}

We use the UCI Heart Disease dataset --- a classic benchmark in clinical ML.

\begin{minted}{python}
# UCI Heart Disease dataset
url = ("https://raw.githubusercontent.com/raphaelfontenelle/"
       "heart-disease-uci/main/heart.csv")
heart = pd.read_csv(url)
print(f"Shape: {heart.shape}")
print(f"Heart disease prevalence: {heart['target'].mean():.1%}")
heart.head()
\end{minted}

\begin{minted}{python}
# Column descriptions
column_info = {
    "age": "Age in years",
    "sex": "1 = male, 0 = female",
    "cp": "Chest pain type (0-3)",
    "trestbps": "Resting blood pressure (mmHg)",
    "chol": "Serum cholesterol (mg/dL)",
    "fbs": "Fasting blood sugar > 120 mg/dL (1 = true)",
    "restecg": "Resting ECG results (0-2)",
    "thalach": "Maximum heart rate achieved",
    "exang": "Exercise-induced angina (1 = yes)",
    "oldpeak": "ST depression induced by exercise",
    "slope": "Slope of peak exercise ST segment",
    "ca": "Number of major vessels colored by fluoroscopy (0-3)",
    "thal": "Thalassemia (1 = normal, 2 = fixed defect, 3 = reversible)",
    "target": "Heart disease (1 = yes, 0 = no)"
}
for col, desc in column_info.items():
    print(f"  {col:12s} : {desc}")
\end{minted}

% ============================================================
\section{Train/test split and cross-validation}

\subsection{Why split the data?}

If you train a model on all your data and then evaluate it on the same data, the accuracy is misleadingly high. The model has ``seen the answers.'' To estimate real-world performance, you must evaluate on data the model has \emph{never seen}.

\begin{minted}{python}
# Define features and target
feature_cols = ["age", "sex", "cp", "trestbps", "chol", "fbs",
                "restecg", "thalach", "exang", "oldpeak", "slope", "ca", "thal"]
X = heart[feature_cols]
y = heart["target"]

# 70/30 train/test split (stratified to preserve class balance)
X_train, X_test, y_train, y_test = train_test_split(
    X, y, test_size=0.3, random_state=42, stratify=y
)
print(f"Training set: {X_train.shape[0]} patients")
print(f"Test set:     {X_test.shape[0]} patients")
print(f"Training prevalence: {y_train.mean():.1%}")
print(f"Test prevalence:     {y_test.mean():.1%}")
\end{minted}

\subsection{Cross-validation}

A single train/test split can be lucky or unlucky. Cross-validation (CV) gives a more robust estimate by rotating the test set:

\begin{minted}{python}
# 5-fold cross-validation with logistic regression
lr = LogisticRegression(max_iter=1000, random_state=42)
cv_scores = cross_val_score(lr, X, y, cv=5, scoring="accuracy")
print(f"CV accuracy: {cv_scores.mean():.3f} +/- {cv_scores.std():.3f}")
print(f"Per-fold: {cv_scores.round(3)}")
\end{minted}

\begin{warningbox}
Never use test set results to choose your model or tune hyperparameters. The test set must only be used \textbf{once}, at the very end. If you peek at test results during model selection, you are effectively training on the test set.
\end{warningbox}

% ============================================================
\section{Decision trees}

A decision tree splits patients into groups using yes/no questions, forming a flowchart that clinicians can follow.

\begin{minted}{python}
# Fit a decision tree
dt = DecisionTreeClassifier(max_depth=4, random_state=42)
dt.fit(X_train, y_train)

# Evaluate
y_pred_dt = dt.predict(X_test)
print("Decision Tree Results:")
print(classification_report(y_test, y_pred_dt,
                            target_names=["No Disease", "Disease"]))
\end{minted}

\begin{minted}{python}
# Visualize the tree
fig, ax = plt.subplots(figsize=(20, 10))
plot_tree(dt, feature_names=feature_cols,
          class_names=["No Disease", "Disease"],
          filled=True, rounded=True, fontsize=9, ax=ax)
plt.title("Decision tree for heart disease prediction")
plt.tight_layout()
plt.show()
\end{minted}

\begin{clinicalnote}
Decision trees are popular in emergency medicine because they produce \emph{clinical decision rules} --- simple flowcharts that a physician can follow at the bedside without a computer. The HEART score, Wells criteria, and Ottawa ankle rules are all conceptually similar to decision trees.
\end{clinicalnote}

\subsection{The danger of deep trees}

\begin{minted}{python}
# Overfit tree (no depth limit)
dt_overfit = DecisionTreeClassifier(random_state=42)  # no max_depth
dt_overfit.fit(X_train, y_train)

print(f"Training accuracy (overfit): {dt_overfit.score(X_train, y_train):.3f}")
print(f"Test accuracy (overfit):     {dt_overfit.score(X_test, y_test):.3f}")
print(f"Training accuracy (depth=4): {dt.score(X_train, y_train):.3f}")
print(f"Test accuracy (depth=4):     {dt.score(X_test, y_test):.3f}")
\end{minted}

\begin{pythontip}
When training accuracy is much higher than test accuracy, your model is \textbf{overfitting} --- it has memorized the training data instead of learning general patterns. Limit tree depth with \texttt{max\_depth} or require a minimum number of samples per leaf with \texttt{min\_samples\_leaf}.
\end{pythontip}

% ============================================================
\section{Random forests}

A random forest combines hundreds of decision trees, each trained on a random subset of the data and features. The final prediction is the majority vote.

\begin{minted}{python}
# Fit a random forest
rf = RandomForestClassifier(n_estimators=200, max_depth=6,
                            random_state=42, n_jobs=-1)
rf.fit(X_train, y_train)

y_pred_rf = rf.predict(X_test)
print("Random Forest Results:")
print(classification_report(y_test, y_pred_rf,
                            target_names=["No Disease", "Disease"]))
\end{minted}

\begin{minted}{python}
# Cross-validation comparison
models = {
    "Logistic Regression": LogisticRegression(max_iter=1000, random_state=42),
    "Decision Tree (depth=4)": DecisionTreeClassifier(max_depth=4, random_state=42),
    "Random Forest": RandomForestClassifier(n_estimators=200, max_depth=6,
                                            random_state=42, n_jobs=-1)
}

print("5-fold CV accuracy:")
for name, model in models.items():
    scores = cross_val_score(model, X, y, cv=5, scoring="accuracy")
    print(f"  {name:30s}: {scores.mean():.3f} +/- {scores.std():.3f}")
\end{minted}

% ============================================================
\section{ROC curves and AUC}

The ROC (Receiver Operating Characteristic) curve plots sensitivity vs (1 -- specificity) across all possible classification thresholds. The AUC (Area Under the Curve) summarizes discriminative ability in a single number.

\begin{itemize}
  \item AUC = 0.5: no better than flipping a coin.
  \item AUC = 0.7--0.8: acceptable discrimination.
  \item AUC = 0.8--0.9: excellent discrimination.
  \item AUC $> 0.9$: outstanding (rare in clinical practice).
\end{itemize}

\begin{minted}{python}
# Get predicted probabilities
lr.fit(X_train, y_train)
y_prob_lr = lr.predict_proba(X_test)[:, 1]
y_prob_rf = rf.predict_proba(X_test)[:, 1]

# Compute ROC curves
fpr_lr, tpr_lr, _ = roc_curve(y_test, y_prob_lr)
fpr_rf, tpr_rf, _ = roc_curve(y_test, y_prob_rf)
auc_lr = roc_auc_score(y_test, y_prob_lr)
auc_rf = roc_auc_score(y_test, y_prob_rf)

# Plot
fig, ax = plt.subplots(figsize=(7, 6))
ax.plot(fpr_lr, tpr_lr, label=f"Logistic Regression (AUC = {auc_lr:.3f})",
        linewidth=2)
ax.plot(fpr_rf, tpr_rf, label=f"Random Forest (AUC = {auc_rf:.3f})",
        linewidth=2)
ax.plot([0, 1], [0, 1], "k--", alpha=0.5, label="Random (AUC = 0.500)")
ax.set_xlabel("False Positive Rate (1 - Specificity)")
ax.set_ylabel("True Positive Rate (Sensitivity)")
ax.set_title("ROC curves: Heart disease prediction")
ax.legend(loc="lower right")
plt.tight_layout()
plt.show()
\end{minted}

\begin{clinicalnote}
In clinical practice, ROC curves are used to compare diagnostic tests. When evaluating a new biomarker for sepsis, you plot its ROC curve against existing biomarkers (CRP, procalcitonin). The test with the highest AUC has the best overall discriminative ability.
\end{clinicalnote}

% ============================================================
\section{Feature importance}

\textbf{Question:} Which risk factors matter most for predicting heart disease?

\begin{minted}{python}
# Random forest feature importance
importances = rf.feature_importances_
importance_df = pd.DataFrame({
    "Feature": feature_cols,
    "Importance": importances
}).sort_values("Importance", ascending=True)

fig, ax = plt.subplots(figsize=(8, 6))
ax.barh(importance_df["Feature"], importance_df["Importance"], color="steelblue")
ax.set_xlabel("Feature importance (Gini)")
ax.set_title("Which risk factors predict heart disease?")
plt.tight_layout()
plt.show()
\end{minted}

\begin{minted}{python}
# Top 5 features
print("Top 5 most important features:")
for _, row in importance_df.tail(5).iloc[::-1].iterrows():
    print(f"  {row['Feature']:12s}: {row['Importance']:.3f}")
\end{minted}

\begin{warningbox}
Feature importance from random forests tells you which variables are useful for \emph{prediction}, not which variables \emph{cause} the outcome. A variable can be important for prediction because it is a proxy for something else. Do not confuse predictive importance with causal importance.
\end{warningbox}

% ============================================================
\section{Calibration}

A model says ``this patient has a 70\% risk of heart disease.'' Does that mean that among all patients the model gives 70\%, roughly 70\% actually have heart disease? If yes, the model is \textbf{well-calibrated}.

\begin{minted}{python}
from sklearn.calibration import calibration_curve

# Calibration plot for random forest
prob_true, prob_pred = calibration_curve(y_test, y_prob_rf, n_bins=8)

fig, ax = plt.subplots(figsize=(6, 6))
ax.plot(prob_pred, prob_true, "o-", linewidth=2, label="Random Forest")
ax.plot([0, 1], [0, 1], "k--", label="Perfectly calibrated")
ax.set_xlabel("Mean predicted probability")
ax.set_ylabel("Observed proportion")
ax.set_title("Calibration plot")
ax.legend()
plt.tight_layout()
plt.show()
\end{minted}

\begin{minted}{python}
# Brier score: lower is better (0 = perfect)
from sklearn.metrics import brier_score_loss

brier_lr = brier_score_loss(y_test, y_prob_lr)
brier_rf = brier_score_loss(y_test, y_prob_rf)
print(f"Brier score (Logistic Regression): {brier_lr:.4f}")
print(f"Brier score (Random Forest):       {brier_rf:.4f}")
\end{minted}

\begin{clinicalnote}
Calibration is critical in clinical prediction models. If a diabetes risk calculator tells a patient ``you have a 40\% chance of developing diabetes in 10 years,'' that number must be accurate --- it drives treatment decisions. The Framingham Risk Score and QRISK are regularly re-calibrated for different populations.
\end{clinicalnote}

% ============================================================
\section{Overfitting: the danger of small clinical datasets}

\begin{minted}{python}
# Demonstrate overfitting with varying dataset sizes
train_sizes = [30, 50, 100, 150, 200, len(X_train)]
results = []

for size in train_sizes:
    if size > len(X_train):
        continue
    X_sub = X_train.iloc[:size]
    y_sub = y_train.iloc[:size]

    rf_temp = RandomForestClassifier(n_estimators=100, random_state=42)
    rf_temp.fit(X_sub, y_sub)

    train_acc = rf_temp.score(X_sub, y_sub)
    test_acc = rf_temp.score(X_test, y_test)
    results.append({"n_train": size, "train_acc": train_acc, "test_acc": test_acc})

results_df = pd.DataFrame(results)

fig, ax = plt.subplots(figsize=(8, 5))
ax.plot(results_df["n_train"], results_df["train_acc"], "o-",
        label="Training accuracy", linewidth=2)
ax.plot(results_df["n_train"], results_df["test_acc"], "s-",
        label="Test accuracy", linewidth=2)
ax.set_xlabel("Number of training patients")
ax.set_ylabel("Accuracy")
ax.set_title("Learning curve: more data reduces overfitting")
ax.legend()
ax.set_ylim(0.5, 1.05)
plt.tight_layout()
plt.show()
\end{minted}

\begin{pythontip}
Signs of overfitting: (1)~training accuracy is much higher than test accuracy, (2)~performance varies wildly across cross-validation folds, (3)~adding more features hurts test performance. The cure: more data, simpler models, regularization, or cross-validation for model selection.
\end{pythontip}

% ============================================================
\section{Practical: heart disease prediction pipeline}

\begin{minted}{python}
# Complete pipeline from raw data to evaluated model
from sklearn.pipeline import Pipeline
from sklearn.impute import SimpleImputer

# Full pipeline
pipeline = Pipeline([
    ("imputer", SimpleImputer(strategy="median")),
    ("scaler", StandardScaler()),
    ("classifier", RandomForestClassifier(n_estimators=200, max_depth=6,
                                          random_state=42))
])

# Cross-validate the full pipeline
cv_scores = cross_val_score(pipeline, X, y, cv=5, scoring="roc_auc")
print(f"Pipeline CV AUC: {cv_scores.mean():.3f} +/- {cv_scores.std():.3f}")

# Final fit and evaluation
pipeline.fit(X_train, y_train)
y_prob_final = pipeline.predict_proba(X_test)[:, 1]
y_pred_final = pipeline.predict(X_test)

print(f"\nFinal test AUC: {roc_auc_score(y_test, y_prob_final):.3f}")
print(f"\n{classification_report(y_test, y_pred_final, target_names=['No Disease', 'Disease'])}")
\end{minted}

% ============================================================
\section{Ethics: fairness in clinical prediction models}

\begin{warningbox}
Clinical prediction models can perpetuate and amplify health disparities. A model trained predominantly on one population may perform poorly on another. Before deploying any clinical ML model, ask:
\begin{itemize}
  \item \textbf{Who is in the training data?} If 90\% of patients are white males, the model may not work for women or minority groups.
  \item \textbf{Does performance differ across subgroups?} Check AUC, sensitivity, and specificity separately by sex, race, and age group.
  \item \textbf{What are the consequences of errors?} A false negative for cardiac risk in a young woman is more dangerous than a false positive.
  \item \textbf{Is the model transparent?} Clinicians and patients deserve to know how a prediction was made.
\end{itemize}
\end{warningbox}

\begin{minted}{python}
# Check model performance by sex
for sex_val, sex_label in [(1, "Male"), (0, "Female")]:
    mask = X_test["sex"] == sex_val
    if mask.sum() < 10:
        print(f"{sex_label}: too few samples ({mask.sum()})")
        continue

    auc_sub = roc_auc_score(y_test[mask], y_prob_final[mask])
    sens = (y_pred_final[mask] == 1)[y_test[mask] == 1].mean()
    spec = (y_pred_final[mask] == 0)[y_test[mask] == 0].mean()

    print(f"{sex_label:8s}: AUC = {auc_sub:.3f}, "
          f"Sensitivity = {sens:.3f}, Specificity = {spec:.3f}, "
          f"n = {mask.sum()}")
\end{minted}

\begin{clinicalnote}
In 2019, a widely-used commercial algorithm for allocating healthcare resources was found to be biased against Black patients. The algorithm used healthcare \emph{costs} as a proxy for healthcare \emph{needs}, but because Black patients had lower costs due to systemic barriers to access, the algorithm systematically underestimated their needs. This affected millions of patients. The lesson: always audit your model for disparities.
\end{clinicalnote}

\begin{exercise}
Using the UCI Heart Disease dataset:
\begin{enumerate}
  \item \textbf{Baseline.} Fit a logistic regression model using all features. Report CV accuracy and AUC.
  \item \textbf{Decision tree.} Fit a decision tree with \texttt{max\_depth=3}. Visualize the tree. Can a clinician follow this decision rule at the bedside?
  \item \textbf{Random forest.} Fit a random forest with 200 trees. Compare its AUC to logistic regression and the decision tree. Is the improvement worth the loss of interpretability?
  \item \textbf{Feature selection.} Train a random forest using only the top 5 most important features. How does performance compare to using all 13 features?
  \item \textbf{ROC comparison.} Plot ROC curves for all three models on the same figure. Which model would you choose for a screening tool? For a confirmatory test?
  \item \textbf{Threshold analysis.} For the random forest, plot sensitivity and specificity as a function of the classification threshold (from 0 to 1). At what threshold are they equal?
  \item \textbf{Calibration.} Plot calibration curves for logistic regression and random forest. Which model is better calibrated?
  \item \textbf{Fairness audit.} Compare AUC, sensitivity, and specificity between male and female patients. Are there concerning disparities? What might explain them?
  \item \textbf{Clinical deployment.} Write a paragraph (in a Markdown cell) describing how you would validate this model before using it in a real hospital. Consider: external validation, regulatory requirements, clinician workflow, and patient consent.
\end{enumerate}
\end{exercise}

% ============================================================
\section{Chapter summary}

\begin{itemize}
  \item Use traditional statistics (regression) for inference and ML for prediction.
  \item Always split data into training and test sets. Never evaluate on training data.
  \item Cross-validation gives more robust performance estimates than a single split.
  \item Decision trees are interpretable and clinician-friendly but prone to overfitting.
  \item Random forests improve accuracy by averaging many trees.
  \item ROC curves and AUC compare the discriminative ability of different models.
  \item Feature importance reveals which risk factors drive predictions (but not causation).
  \item Calibration checks whether predicted probabilities match observed outcomes.
  \item Overfitting is the central danger of ML on small clinical datasets --- always validate.
  \item Fairness audits are essential: check model performance across demographic subgroups before deployment.
\end{itemize}
